\documentclass[dvipdfmx,11pt,a4paper]{article}
\usepackage{amsmath,amssymb,amsthm,amsfonts}
\usepackage{geometry}
\geometry{margin=1.1in}
\usepackage{enumitem}
\usepackage{booktabs}
\usepackage{mathtools}
\usepackage[dvipdfmx]{hyperref}
\usepackage[nameinlink,noabbrev]{cleveref}

\newtheorem{theorem}{Theorem}[section]
\newtheorem{proposition}[theorem]{Proposition}
\newtheorem{lemma}[theorem]{Lemma}
\newtheorem{corollary}[theorem]{Corollary}
\theoremstyle{definition}
\newtheorem{definition}[theorem]{Definition}
\theoremstyle{remark}
\newtheorem{remark}[theorem]{Remark}

\newcommand{\Jnorm}{J_{\rm norm}}

\title{D4L and Hilbert Space:\\
A Sibling Construction, not an Identification}

\author{https://github.com/hypernumbernet}
\date{\today}

\begin{document}
\maketitle

\begin{abstract}
Dual Spacetime 4-Valued Logic (D4L) is not Birkhoff--von Neumann quantum
logic. The parameter Killing form is indefinite of signature $(3,3)$;
self-overlap can be negative and is therefore not a Born probability;
the scalar connectives $\min/\max$ are distributive. What DST does
supply is a genuine Hilbert space of dual-sector kinematics: the cyclic
generators satisfy the quaternion table, the flipped dual pairing is
Euclidean on $\mathbb{R}^3$, and $-i\sigma_a$ represents those generators
on $\mathbb{C}^2$. The closed-subspace lattice of $\mathbb{C}^2$ is
orthomodular and non-distributive. Four D4L labels cannot be four
pairwise orthogonal rays in that space. All numbered theorems below are
machine-checked in Lean~4. None of this identifies D4L with an
orthomodular lattice, and none of it derives the Born rule from $J$.
\end{abstract}

\section{Position}
\label{sec:pos}

D4L \cite{D4L2026} takes admissible torsion configurations as
propositions, $\Jnorm$ as the observable, and $\min/\max/-$ as scalar
connectives. Its geometric layer uses the parameter Killing form and
the bivector commutator \cite{Diophantine2026}. Those objects were
deliberately \emph{not} called a Hilbert inner product, a Born rule, or
an orthomodular lattice.

This note constructs the Hilbert space that the dual sector already
contains, and records the separations that forbid reading D4L as
Birkhoff--von Neumann logic. The construction is a sibling of D4L, not
a rebranding. The Lean entry remains $\texttt{DstDiophantine.Logic}$;
the new modules live under $\texttt{Logic.Quantum}$ and are not
re-exported from $\texttt{DstDiophantine.Basic}$.

The identification $I\leftrightarrow i\sigma$ appearing in older
quantum-gravity drafts is replaced by $I\mapsto -i\sigma$. That sign
makes $IJ=K$ agree with the cyclic products in $G(3,1,1)$ and sends the
dual rotor to the standard $\mathrm{SU}(2)$ factor
$\exp(-i\beta\cdot\sigma/2)$. Lean is authoritative.

\section{Separations}
\label{sec:sep}

Write $p=(\alpha_a,\beta_a)_{a=0,1,2}$ and
\[
\texttt{killingForm}(p,q)
= 8\sum_{a=0}^{2}\bigl(\alpha_a\alpha'_a-\beta_a\beta'_a\bigr).
\]
Self-overlap is $\langle p\mid p\rangle=16\,J(p)$.

\begin{theorem}[Indefinite pairing]
\label{thm:indef}
There exist configurations with $\texttt{killingForm}(p,p)>0$ and with
$\texttt{killingForm}(q,q)<0$. In particular the pairing is not
positive definite. Lean: $\texttt{killingForm\_indefinite}$.
\end{theorem}

The witnesses are the pure-boost and pure-rotation extremals already
used for $\lvert\Jnorm\rvert=1$.

\begin{theorem}[Not a Born probability]
\label{thm:born}
There exists $p$ with $\langle p\mid p\rangle<0$. A signed quantity
cannot be a Born probability. Lean:
$\texttt{overlap\_self\_not\_born\_probability}$.
\end{theorem}

The identification $\lvert\psi\rvert^2\propto B(\Omega,\Omega)$ from
the quantum-gravity companion is therefore not formalised here.

\begin{theorem}[Distributive scalars]
\label{thm:dist}
For all $a,b,c\in\mathbb{R}$,
$a\wedge(b\vee c)=(a\wedge b)\vee(a\wedge c)$ with
$\wedge=\min$ and $\vee=\max$. Lean:
$\texttt{scalar\_connectives\_distributive}$.
\end{theorem}

A distributive lattice cannot be isomorphic to the subspace lattice of
a Hilbert space of dimension at least two.

\section{Dual and usual sectors}
\label{sec:sector}

\begin{definition}
$p$ is dual-sector when $\alpha=0$, usual-sector when $\beta=0$.
Write $\texttt{ofDual}\,\beta$ and $\texttt{ofUsual}\,\alpha$ for the
corresponding embeddings of Euclidean $\mathbb{R}^3$.
\end{definition}

\begin{theorem}[Definite restrictions]
\label{thm:sector}
On the dual sector,
$\texttt{killingForm}(\texttt{ofDual}\,\beta,\texttt{ofDual}\,\gamma)
=-8\langle\beta,\gamma\rangle_{\mathbb{R}^3}$,
hence negative definite. On the usual sector the pairing is
$+8\langle\alpha,\gamma\rangle_{\mathbb{R}^3}$, hence positive
definite. Lean: $\texttt{killingForm\_ofDual}$,
$\texttt{killingForm\_ofDual\_neg\_def}$.
\end{theorem}

The inner-product space is the linear span of each sector, not the
admissible cube $\alpha_a,\beta_a\ge 0$, $\alpha_a+\beta_a\le\pi/2$.
Superposition is a statement about that span.

\section{Quaternion table}
\label{sec:quat}

The cyclic generators already satisfy
$(\texttt{cyclic}\,a)^2=-1$. The off-diagonal products are new.

\begin{theorem}[Quaternion multiplication]
\label{thm:quat}
\[
\begin{aligned}
\texttt{cyclic}\,0\cdot\texttt{cyclic}\,1&=\texttt{cyclic}\,2,\\
\texttt{cyclic}\,1\cdot\texttt{cyclic}\,2&=\texttt{cyclic}\,0,\\
\texttt{cyclic}\,2\cdot\texttt{cyclic}\,0&=\texttt{cyclic}\,1,
\end{aligned}
\qquad
\begin{aligned}
\texttt{cyclic}\,1\cdot\texttt{cyclic}\,0&=-\texttt{cyclic}\,2,\\
\texttt{cyclic}\,2\cdot\texttt{cyclic}\,1&=-\texttt{cyclic}\,0,\\
\texttt{cyclic}\,0\cdot\texttt{cyclic}\,2&=-\texttt{cyclic}\,1.
\end{aligned}
\]
Lean: $\texttt{cyclic\_zero\_mul\_one}$ and companions.
\end{theorem}

The identities are computed from the vector anticommutators in
$G(3,1,1)$. For a pure dual configuration one has
$\Omega(\texttt{ofDual}\,\beta)=\sum(\beta_a/2)\,\texttt{cyclic}\,a$,
and on axis $0$ this collapses to $(\theta/2)\,\texttt{cyclic}\,0$.

\section{Spinor Hilbert space}
\label{sec:spinor}

\begin{definition}[Dual spinor]
$\texttt{DualSpinor}:=\mathbb{C}^2$ with the standard Hermitian
product. Pauli matrices $\sigma_x,\sigma_y,\sigma_z$ are the usual
ones. The representation is
$\texttt{cyclicRep}\,a:=-i\sigma_a$.
\end{definition}

\begin{theorem}[Representation]
\label{thm:rep}
$(\texttt{cyclicRep}\,a)^2=-1$ and
$\texttt{cyclicRep}\,0\cdot\texttt{cyclicRep}\,1=\texttt{cyclicRep}\,2$.
Each $\texttt{cyclicRep}\,a$ is anti-Hermitian. Lean:
$\texttt{cyclicRep\_sq}$, $\texttt{cyclicRep\_zero\_mul\_one}$,
$\texttt{cyclicRep\_conjTranspose}$.
\end{theorem}

The dual rotor is the matrix exponential
\[
\texttt{dualRotorMat}(\beta)
=\exp\Bigl(\sum_a(\beta_a/2)(-i\sigma_a)\Bigr)
=\exp(-i\beta\cdot\sigma/2).
\]
On a single axis this is exactly $\exp$ of the corresponding generator.
A general identification of the Banach exponential in $G(3,1,1)$ with
the matrix exponential is not claimed; that would require BCH control
already refused by $\texttt{Algebra.Motor}$.

$\dim_{\mathbb{C}}\texttt{DualSpinor}=2$, so a projective measurement
has at most two orthogonal outcomes.

\section{Quantum logic of $\mathbb{C}^2$}
\label{sec:ql}

\begin{definition}
A quantum proposition is a linear subspace of $\texttt{DualSpinor}$.
Meet is intersection, join is the sum, negation is the orthogonal
complement. Finite dimension makes every subspace closed.
\end{definition}

\begin{theorem}[Orthomodular law]
\label{thm:om}
If $A\le B$ then $A\vee(A^\perp\wedge B)=B$. Lean:
$\texttt{orthomodular}$.
\end{theorem}

This is the orthogonal decomposition of a finite-dimensional subspace,
already in mathlib.

\begin{theorem}[Non-distributivity]
\label{thm:ndist}
Let $e_0=(1,0)$, $e_1=(0,1)$, $d=(1,1)$, and write $\ell_v$ for
$\mathbb{C}v$. Then
\[
\ell_{e_0}\wedge(\ell_{e_1}\vee\ell_d)
\neq
(\ell_{e_0}\wedge\ell_{e_1})\vee(\ell_{e_0}\wedge\ell_d).
\]
The left side is $\ell_{e_0}$; the right side is $\{0\}$. Lean:
$\texttt{not\_distributive}$.
\end{theorem}

\begin{theorem}[At most two orthogonal atoms]
\label{thm:two}
There do not exist three pairwise orthogonal one-dimensional subspaces
of $\texttt{DualSpinor}$. Lean:
$\texttt{no\_three\_pairwise\_orthogonal\_lines}$.
\end{theorem}

\section{Dictionary}
\label{sec:dict}

\begin{center}
\begin{tabular}{@{}ll@{}}
\toprule
D4L & Dual Hilbert / quantum logic \\
\midrule
pure dual amplitude & $\beta\in\mathbb{R}^3$ and $U(\beta)\in\mathrm{U}(2)$ \\
usual--dual swap & leaves the dual sector; not a Hilbert adjoint \\
Killing overlap & distinct from the spinor inner product \\
$\texttt{interfere}=0$ on one axis & commuting generators (necessary only) \\
$\Jnorm$ four-state collapse & coarse-graining of a continuous scalar \\
$\min/\max$ & not the meet/join of $L(\mathbb{C}^2)$ \\
\bottomrule
\end{tabular}
\end{center}

\begin{theorem}[Four labels are not a PVM]
\label{thm:pvm}
There is no assignment of one-dimensional subspaces of
$\texttt{DualSpinor}$ to the four D4L labels that makes distinct
labels pairwise orthogonal. Lean:
$\texttt{four\_labels\_not\_orthogonal\_pvm}$.
\end{theorem}

\section{What this note does and does not show}
\label{sec:claims}

\begin{itemize}[leftmargin=1.4em]
  \item It \emph{does} prove that D4L's pairing and connectives cannot
        be a Hilbert inner product, a Born rule, or an orthomodular
        lattice.
  \item It \emph{does} extract a Euclidean dual sector and a quaternion
        table for the cyclic generators.
  \item It \emph{does} represent that table on the Hilbert space
        $\mathbb{C}^2$ and record the orthomodular, non-distributive
        lattice of its subspaces.
  \item It does \emph{not} identify D4L with Birkhoff--von Neumann
        logic.
  \item It does \emph{not} derive the Born rule from $J$ or from
        Killing overlap.
  \item It does \emph{not} construct the Dirac spinor $\mathbb{C}^4$,
        position/momentum operators, or a solution of the measurement
        problem.
  \item It does \emph{not} identify the D4L scale $\lambda$ with
        Hilbert dimension, CGA dilation, or $\texttt{scaleTorsion}$.
\end{itemize}

The Lean modules are $\texttt{Separation}$, $\texttt{DualSector}$,
$\texttt{Quaternion}$, $\texttt{Spinor}$, $\texttt{QuantumLogic}$, and
$\texttt{Dictionary}$, under $\texttt{DstDiophantine.Logic.Quantum}$.

\begin{thebibliography}{9}
\raggedright

\bibitem{D4L2026}
Dual Spacetime 4-Valued Logic (D4L): Amplitudes, Four States, and
Geometric Interference,
\texttt{dst-4-valued-logic.tex} (2026).

\bibitem{Diophantine2026}
Discrete Biquaternionic Double Spacetime,
\texttt{dst-diophantine.tex} (2026). Lean companion:
\url{https://github.com/hypernumbernet/DstDiophantine}.

\bibitem{BvN1936}
G.~Birkhoff and J.~von Neumann, The logic of quantum mechanics,
Ann.\ Math.\ 37 (1936).

\end{thebibliography}

\end{document}
